% =========================================================
% TIG3 MAIN MANUSCRIPT
% Topological Integrity Gravity III
% =========================================================

\documentclass[11pt]{article}

% =========================================================
% Packages
% =========================================================

\usepackage[a4paper,margin=1in]{geometry}

\usepackage{amsmath}
\usepackage{amssymb}
\usepackage{amsthm}
\usepackage{bm}
\usepackage{mathrsfs}

\usepackage{graphicx}
\usepackage{float}

\usepackage{hyperref}
\usepackage{microtype}

\usepackage{enumitem}

% =========================================================
% Hyperref
% =========================================================

\hypersetup{
    colorlinks=true,
    linkcolor=blue,
    citecolor=blue,
    urlcolor=blue
}

% =========================================================
% Theorem Environments
% =========================================================

\newtheorem{theorem}{Theorem}
\newtheorem{proposition}{Proposition}
\newtheorem{remark}{Remark}

% =========================================================
% Title
% =========================================================

\title{
\Large
Topological Integrity Gravity III:\\[0.3cm]
\large
Conditional Vacuum Admissibility and\\
Bounded Curvature Structure
}

\author{
Kai Stefan Dietrich\\
Independent Researcher
}

\date{\today}

% =========================================================
% Document
% =========================================================

\begin{document}

\maketitle

% =========================================================
% Abstract
% =========================================================

\begin{abstract}

We develop a conditional operator-theoretic perturbative framework for a representative bounded-curvature vacuum sector compatible with the cubic branch structure established in Topological Integrity Gravity II (TIG2).

The framework is based on a structurally regulated spherically symmetric geometry preserving Schwarzschild-compatible asymptotics while introducing a finite representative core scale. We analyze the resulting perturbative radial operator structure and formulate conditions for quadratic-form control, self-adjoint realization, spectral admissibility, and bounded linear perturbative evolution.

The present work is formulated conservatively as a perturbative admissibility framework rather than a complete nonlinear gravitational theory. In particular, the analysis does not establish nonlinear stability, complete covariant field equations, quantum-gravity completion, or full singularity resolution.

The objective of the present paper is to define a mathematically coherent vacuum-sector program suitable for further operator-theoretic and geometric analysis within a publication-stable mathematical physics framework.

\end{abstract}

% =========================================================
% Table of Contents
% =========================================================

\tableofcontents

\newpage

% =========================================================
% Sections
% =========================================================

% =========================================================
% Introduction and Geometry
% =========================================================

\section{Introduction}

\subsection{Motivation}

One of the central open problems in gravitational physics concerns the structural behavior of spacetime geometry in high-curvature vacuum regimes. In particular, classical Schwarzschild geometry predicts the formation of curvature singularities associated with geodesic incompleteness and divergent curvature invariants \cite{wald1984,hawkingellis1973,penrose1965,hawking1970}.

A variety of approaches have attempted to regulate such singular structures through modified geometric sectors, effective matter models, or quantum-gravitational corrections \cite{bardeen1968,dymnikova1992,hayward2006,ansoldi2008}. However, many such constructions either introduce uncontrolled ultraviolet assumptions or lack mathematically stable perturbative formulations.

The present work develops a conservative operator-theoretic perturbative framework associated with a representative bounded-curvature vacuum geometry compatible with the branch structure established in Topological Integrity Gravity II (TIG2).

The goal of the present paper is not to propose a complete alternative theory of gravity, but rather to define a mathematically coherent perturbative admissibility framework suitable for further spectral and geometric analysis.

% =========================================================

\subsection{TIG2 Structural Background}

The present framework builds upon the cubic branch structure introduced in TIG2 \cite{tig2}.

The central structural relation is

\begin{equation}
x^3 - x^2 + \beta^3 = 0,
\end{equation}

where

\begin{equation}
x = \frac{r_H}{2M},
\qquad
\beta = \frac{r_c}{2M}.
\end{equation}

Here:
\begin{itemize}[leftmargin=0.6cm]
\item $r_H$ denotes the horizon radius,
\item $M$ is the mass parameter,
\item and $r_c$ is a regulated representative core scale.
\end{itemize}

The discriminant structure is given by

\begin{equation}
\Delta(\beta)
=
\beta^3(4 - 27\beta^3).
\end{equation}

The critical branch-merger point occurs at

\begin{equation}
x_c = \frac{2}{3},
\qquad
\beta_c^3 = \frac{4}{27}.
\end{equation}

This branch structure defines the admissible near-critical sector underlying the TIG3 vacuum program.

% =========================================================

\subsection{Scientific Scope}

The present paper develops a conditional perturbative operator framework associated with a representative bounded-curvature geometry.

The analysis focuses on:
\begin{itemize}[leftmargin=0.6cm]
\item representative vacuum geometry,
\item perturbative admissibility,
\item operator-theoretic structure,
\item spectral consistency,
\item and bounded linear evolution.
\end{itemize}

The framework does not currently establish:
\begin{itemize}[leftmargin=0.6cm]
\item nonlinear stability,
\item complete covariant field equations,
\item quantum-gravity completion,
\item observational confirmation,
\item or full singularity resolution.
\end{itemize}

All perturbative statements are conditional upon the operator-theoretic assumptions explicitly stated below.

% =========================================================

\section{Representative TIG3 Geometry}

\subsection{Metric Ansatz}

We consider a static spherically symmetric representative geometry of the form

\begin{equation}
ds^2
=
-F(r)\,dt^2
+
\frac{dr^2}{F(r)}
+
r^2 d\Omega^2,
\end{equation}

where

\begin{equation}
d\Omega^2
=
d\theta^2
+
\sin^2\theta\, d\phi^2.
\end{equation}

The representative lapse function is chosen as

\begin{equation}
F(r)
=
1
-
\frac{2Mr^2}{r^3 + r_c^3}.
\label{eq:lapse}
\end{equation}

The geometry preserves Schwarzschild-compatible asymptotics while introducing a regulated high-curvature core scale.

% =========================================================

\subsection{Asymptotic Structure}

For large radial distance,

\begin{equation}
r \gg r_c,
\end{equation}

the lapse function admits the asymptotic expansion

\begin{equation}
F(r)
=
1
-
\frac{2M}{r}
+
\mathcal{O}\!\left(\frac{r_c^3}{r^3}\right).
\end{equation}

Hence the representative geometry reproduces Schwarzschild-compatible asymptotic behavior in the low-curvature regime \cite{wald1984,chandrasekhar1983}.

This compatibility condition is essential for maintaining consistency with the classical far-field sector of General Relativity.

% =========================================================

\subsection{Near-Core Structure}

Near the regulated core,

\begin{equation}
r \ll r_c,
\end{equation}

the lapse function behaves as

\begin{equation}
F(r)
\approx
1
-
\frac{2M}{r_c^3}r^2.
\end{equation}

The resulting geometry possesses bounded representative curvature invariants within the analyzed sector.

In particular, the representative Kretschmann scalar remains finite near the regulated core.

The present work interprets this structure conservatively as a bounded-curvature representative geometry rather than a complete singularity-resolution theorem.

% =========================================================

\subsection{Horizon Structure}

The horizon sector is determined by the roots of

\begin{equation}
F(r_H)=0.
\end{equation}

After rescaling according to

\begin{equation}
x
=
\frac{r_H}{2M},
\qquad
\beta
=
\frac{r_c}{2M},
\end{equation}

the horizon equation reduces exactly to the TIG2 cubic branch relation

\begin{equation}
x^3 - x^2 + \beta^3 = 0.
\end{equation}

Thus the representative TIG3 geometry is structurally compatible with the admissible branch hierarchy established in TIG2 \cite{tig2}.

% =========================================================

\subsection{Near-Critical Sector}

Near the critical branch-merger regime,

\begin{equation}
\beta \approx \beta_c,
\end{equation}

the geometry may develop elongated throat structures and modified perturbative trapping behavior.

The present work does not establish a rigorous near-critical resonance theorem. Such structures are interpreted only as conjectural perturbative sectors requiring further operator-theoretic analysis.

% =========================================================

\subsection{Interpretation}

The representative geometry introduced above should be interpreted as a mathematically controlled bounded-curvature sector suitable for perturbative operator analysis.

The present framework is intentionally conservative.

In particular, the geometry is not currently claimed to represent:
\begin{itemize}[leftmargin=0.6cm]
\item a complete covariant gravitational theory,
\item a final black-hole replacement model,
\item or a quantum-gravitational completion.
\end{itemize}

Rather, the geometry serves as the background structure for the perturbative admissibility program developed in the following sections.


% =========================================================
% Perturbative Operator Framework
% =========================================================

\section{Perturbative Operator Framework}

\subsection{Perturbative Structure}

We consider perturbations of the representative background geometry introduced in the previous section.

The metric perturbation is written as

\begin{equation}
g_{ab}
\rightarrow
g_{ab} + h_{ab},
\end{equation}

where the perturbative sector satisfies

\begin{equation}
|h_{ab}| \ll 1.
\end{equation}

The present analysis is restricted to the linear perturbative regime.

No nonlinear stability claims are made.

% =========================================================

\subsection{Spherical Symmetry and Mode Reduction}

Due to the static spherical symmetry of the representative background geometry, perturbative modes may be decomposed into angular sectors.

The perturbative analysis is formulated using a reduced radial description associated with angular harmonic decomposition.

The present work does not attempt a complete tensor-harmonic classification. Instead, we focus on the resulting effective radial perturbative structure.

This reduced framework is interpreted as a representative operator-theoretic sector suitable for perturbative admissibility analysis.

% =========================================================

\subsection{Tortoise Coordinate}

We introduce the tortoise coordinate $r_*$ through

\begin{equation}
\frac{dr_*}{dr}
=
\frac{1}{F(r)},
\end{equation}

where $F(r)$ is the representative lapse function defined in Eq.~(\ref{eq:lapse}).

The tortoise coordinate regularizes the perturbative radial structure near simple horizon sectors and allows the perturbation equations to be expressed in Schrödinger-type form.

% =========================================================

\subsection{Effective Radial Operator}

The perturbative radial dynamics are modeled formally through the operator

\begin{equation}
L_{\mathrm{TIG}}
=
-
\frac{d^2}{dr_*^2}
+
V_{\mathrm{eff}}(r).
\label{eq:operator}
\end{equation}

The effective potential is represented by the candidate structure

\begin{equation}
V_{\mathrm{eff}}(r)
=
F(r)
\left[
\frac{\ell(\ell+1)}{r^2}
+
\frac{F'(r)}{r}
\right].
\label{eq:potential}
\end{equation}

Here $\ell$ denotes the angular harmonic index.

The exact perturbative derivation of the effective potential remains incomplete and is treated conservatively throughout the present work.

Accordingly, Eq.~(\ref{eq:potential}) should be interpreted as a representative candidate perturbative potential compatible with the analyzed geometric sector.

% =========================================================

\subsection{Asymptotic Behavior of the Potential}

For large radial distance,

\begin{equation}
r \gg r_c,
\end{equation}

the effective potential approaches the Schwarzschild-type asymptotic structure

\begin{equation}
V_{\mathrm{eff}}(r)
\sim
\frac{\ell(\ell+1)}{r^2}
+
\mathcal{O}\!\left(\frac{1}{r^3}\right).
\end{equation}

Thus the perturbative sector remains compatible with the classical far-field regime.

Near the regulated core, the bounded behavior of $F(r)$ prevents the emergence of divergent representative potential terms within the analyzed perturbative sector.

% =========================================================

\subsection{Quadratic Form Structure}

Associated with the formal operator (\ref{eq:operator}), we define the quadratic form

\begin{equation}
Q[\psi]
=
\int
\left(
\left|
\frac{d\psi}{dr_*}
\right|^2
+
V_{\mathrm{eff}}(r)
|\psi|^2
\right)
dr_*.
\label{eq:quadratic_form}
\end{equation}

The quadratic form plays a central role in the perturbative admissibility structure.

In particular, lower-boundedness of $Q[\psi]$ is required for controlled spectral behavior and bounded linear evolution.

The present work does not establish a complete positivity theorem for the quadratic form.

Such results remain conditional upon further operator-theoretic analysis.

% =========================================================

\subsection{Self-Adjoint Realization}

The perturbative framework requires a mathematically controlled realization of the formal operator $L_{\mathrm{TIG}}$.

The present analysis adopts a conservative operator-theoretic perspective in which self-adjoint realization is treated conditionally.

Formally, one seeks a self-adjoint extension

\begin{equation}
L_{\mathrm{TIG}}^{F},
\end{equation}

obtained through quadratic-form methods and admissible endpoint classification.

The existence and uniqueness of such a realization depend on:
\begin{itemize}[leftmargin=0.6cm]
\item endpoint behavior,
\item domain structure,
\item quadratic-form closability,
\item and lower-boundedness conditions.
\end{itemize}

The complete proof of self-adjoint realization remains open.

% =========================================================

\subsection{Spectral Admissibility}

The central spectral admissibility condition of the present framework is

\begin{equation}
\sigma(L_{\mathrm{TIG}}^{F})
\subset
[0,\infty),
\label{eq:spectral_condition}
\end{equation}

where $\sigma(L_{\mathrm{TIG}}^{F})$ denotes the spectrum of the self-adjoint perturbative operator.

Condition (\ref{eq:spectral_condition}) excludes negative spectral modes corresponding to exponentially growing perturbative instabilities.

The present work does not establish a rigorous spectral nonnegativity theorem.

Instead, Eq.~(\ref{eq:spectral_condition}) defines the conditional spectral admissibility requirement underlying the TIG3 perturbative program.

% =========================================================

\subsection{Conditional Bounded Linear Evolution}

Under the assumptions:
\begin{itemize}[leftmargin=0.6cm]
\item admissible endpoint structure,
\item self-adjoint realization,
\item and spectral nonnegativity,
\end{itemize}

the perturbative evolution equation

\begin{equation}
\partial_t^2 \Psi
+
L_{\mathrm{TIG}}^{F}\Psi
=
0
\end{equation}

admits bounded linear finite-energy evolution within the analyzed perturbative sector.

This statement is conditional.

No nonlinear boundedness theorem is established in the present work.

% =========================================================

\subsection{Near-Critical Perturbative Sector}

Near the critical branch-merger regime

\begin{equation}
\beta \approx \beta_c,
\end{equation}

the perturbative geometry may develop modified spectral structures associated with elongated throat behavior.

Potential effects include:
\begin{itemize}[leftmargin=0.6cm]
\item slow decay,
\item metastable trapping,
\item resonance accumulation,
\item and spectral compression.
\end{itemize}

The present work does not establish rigorous resonance theorems in the near-critical sector.

Such structures remain conjectural and require future asymptotic spectral analysis.

% =========================================================

\subsection{Interpretation}

The perturbative framework developed above should be interpreted as a conditional operator-theoretic admissibility program associated with a representative bounded-curvature vacuum geometry.

The present analysis establishes:
\begin{itemize}[leftmargin=0.6cm]
\item a coherent perturbative radial operator structure,
\item a quadratic-form framework,
\item and a conditional spectral admissibility hierarchy.
\end{itemize}

The framework does not currently establish:
\begin{itemize}[leftmargin=0.6cm]
\item nonlinear stability,
\item complete spectral classification,
\item exact perturbative derivation of the effective potential,
\item or full covariant gravitational closure.
\end{itemize}

The present work is therefore interpreted conservatively as a mathematically structured perturbative admissibility framework suitable for further operator-theoretic development.


% =========================================================
% Conclusions and Limitations
% =========================================================

\section{Conclusions and Limitations}

\subsection{Summary of Results}

The present work develops a conditional operator-theoretic perturbative framework associated with a representative bounded-curvature vacuum geometry compatible with the branch structure established in TIG2 \cite{tig2}.

The framework is constructed around the representative lapse function

\begin{equation}
F(r)
=
1
-
\frac{2Mr^2}{r^3+r_c^3},
\end{equation}

which preserves Schwarzschild-compatible asymptotic behavior while introducing a finite representative core scale \cite{wald1984,chandrasekhar1983}.

The horizon sector reduces exactly to the cubic TIG2 branch relation

\begin{equation}
x^3 - x^2 + \beta^3 = 0,
\end{equation}

with critical branch-merger structure

\begin{equation}
x_c = \frac{2}{3},
\qquad
\beta_c^3 = \frac{4}{27}.
\end{equation}

The perturbative dynamics are modeled through a reduced radial operator framework associated with a candidate effective perturbative potential.

The resulting construction yields:
\begin{itemize}[leftmargin=0.6cm]
\item a representative bounded-curvature geometry,
\item a perturbative radial operator structure,
\item a quadratic-form framework,
\item a conditional self-adjointness program,
\item and a conditional spectral admissibility hierarchy.
\end{itemize}

% =========================================================

\subsection{Interpretation of the Present Framework}

The present work should be interpreted conservatively as a structured perturbative admissibility program within mathematical physics.

The framework does not attempt to establish a complete gravitational theory.

Rather, the objective is to define a mathematically coherent vacuum-sector structure suitable for further operator-theoretic and geometric analysis.

In particular, the present analysis develops:
\begin{itemize}[leftmargin=0.6cm]
\item perturbative consistency structure,
\item admissibility conditions,
\item and conditional bounded linear evolution.
\end{itemize}

The framework therefore represents a mathematically controlled perturbative sector rather than a final gravitational completion.

% =========================================================

\subsection{Conditional Nature of the Results}

Several central components of the framework remain conditional.

In particular:
\begin{itemize}[leftmargin=0.6cm]
\item the exact derivation of the effective perturbative potential,
\item endpoint classification,
\item quadratic-form positivity,
\item and full self-adjoint realization
\end{itemize}

remain open operator-theoretic problems \cite{friedrichs1934,reedsimon2,reedsimon4,kato1995,weidmann1987}.

Accordingly, the spectral admissibility condition

\begin{equation}
\sigma(L_{\mathrm{TIG}}^{F})
\subset
[0,\infty)
\end{equation}

is presently interpreted as a conditional admissibility requirement rather than a rigorously established theorem.

% =========================================================

\subsection{Non-Claims}

The present work does not establish:
\begin{itemize}[leftmargin=0.6cm]
\item nonlinear stability,
\item complete spectral classification,
\item full covariant field equations,
\item quantum-gravity completion,
\item observational verification,
\item or complete singularity-resolution theorems.
\end{itemize}

The framework also does not claim to replace General Relativity.

Instead, the representative geometry is constructed specifically to preserve Schwarzschild-compatible asymptotics within the low-curvature regime \cite{wald1984,hawkingellis1973}.

% =========================================================

\subsection{Near-Critical Sector}

Near the critical branch-merger regime

\begin{equation}
\beta \approx \beta_c,
\end{equation}

the representative geometry may develop modified perturbative spectral structure associated with elongated throat behavior and metastable trapping sectors.

Potential effects include:
\begin{itemize}[leftmargin=0.6cm]
\item resonance accumulation,
\item slow decay,
\item spectral compression,
\item and long-lived perturbative sectors.
\end{itemize}

The present work does not establish rigorous resonance theorems in this regime.

Such structures remain conjectural and require separate asymptotic spectral analysis.

% =========================================================

\subsection{Scientific Position}

The strongest currently justified interpretation of the present framework is the following:

\begin{quote}
TIG3 provides a coherent conditional operator-theoretic perturbative admissibility framework for a representative bounded-curvature vacuum sector compatible with the TIG2 cubic branch structure.
\end{quote}

This statement defines the intended scientific scope of the present work.

% =========================================================

\subsection{Future Directions}

Several mathematically significant problems remain open.

These include:
\begin{itemize}[leftmargin=0.6cm]
\item rigorous derivation of the effective perturbative potential,
\item endpoint classification,
\item quadratic-form coercivity,
\item spectral nonnegativity,
\item resonance analysis in the near-critical regime,
\item and possible nonlinear extensions of the perturbative framework.
\end{itemize}

Future work may also investigate whether the present perturbative admissibility structure can be embedded into a more complete covariant geometric formulation.

Such developments remain outside the scope of the present paper.

% =========================================================

\subsection{Final Remarks}

The present work establishes a mathematically conservative perturbative admissibility framework associated with a representative bounded-curvature vacuum geometry compatible with the TIG2 branch structure.

The framework is intended as a structurally controlled mathematical physics program emphasizing:
\begin{itemize}[leftmargin=0.6cm]
\item explicit assumptions,
\item theorem transparency,
\item operator-theoretic consistency,
\item and publication-stable scientific interpretation.
\end{itemize}

The present analysis therefore represents a structured continuation program in perturbative geometric operator theory rather than a final theory of gravitation.


% =========================================================
% Bibliography
% =========================================================

\bibliographystyle{plain}

\bibliography{references}

% =========================================================
% End
% =========================================================

\end{document}
